% !TEX program = pdflatex
% LLMWorld - full paper (v2). Single-column article; pdfLaTeX-compatible, all-ASCII.
\documentclass[11pt]{article}

\usepackage[utf8]{inputenc}
\usepackage[T1]{fontenc}
\usepackage{lmodern}
\usepackage[a4paper,margin=2.4cm]{geometry}
\usepackage{amsmath}
\usepackage{amssymb}
\usepackage{array}
\usepackage{graphicx}
\usepackage{booktabs}
\usepackage{enumitem}
\usepackage[hidelinks]{hyperref}

\setlist{nosep,leftmargin=1.5em}

\title{\bfseries From Language to Load:\\ How Household LLM Agents Respond to Pricing, Norms and Disruptions\\ --- and Where the Response Stops Being Real}
\author{Hongkun Lyu (Student ID 35933771)\\ Supervisor: Hao Wang\\ \small FIT5216 --- final paper}
\date{September 2026}

\begin{document}
\maketitle

%======================================================================
\begin{abstract}
\noindent
Residential electricity demand is increasingly shaped by heterogeneous household behaviour and by
distributed resources, yet the two dominant modelling traditions remain limited: top-down forecasting
treats households as an undifferentiated mass, and classical bottom-up models encode behaviour as fixed
states and transition probabilities that cannot process rich, contextual or textual signals such as
dynamic tariffs, social norms or news of a heatwave. Large language model (LLM) generative agents offer a
far more expressive behavioural engine, but existing LLM-based energy simulators are almost exclusively
single-household or small-scale, and none has been systematically validated as a \emph{population-level
policy-intervention} testbed against the empirical demand-response literature. This paper presents
\textsc{LLMWorld}, an end-to-end multi-agent framework in which each household is an autonomous LLM agent
whose day is planned from a structured profile and then mapped to minute-resolution, appliance-level load;
crucially, \emph{every} external input --- time-of-use or critical-peak prices, taxes, rebates, social
norms, and news events --- is delivered \emph{only in natural language}, so that aggregate behaviour is
emergent rather than encoded. We run a large, controlled battery of experiments on three heterogeneous
synthetic communities under four validity rules derived from failure (same-era baselines, $n\ge15$,
cross-world replication, and content-contrast), and we deliberately subject every conclusion to an
adversarial audit. We obtain seven statistically significant cross-world event results, including a
previously unreported \emph{stay-home family} --- public holidays, work-from-home days and transport
strikes --- that each raise a household's daily energy by roughly one fifth to two thirds, and we
characterise the evening peak as cooking-dominated and therefore largely resistant to generic signals. We
further show that time-of-use and critical-peak prices, which are \emph{null} as plain text, become
significantly peak-shifting once the trade-off is made concrete through an injected per-appliance
peak/off-peak cost table; decomposing this effect, we find that the money figures alone already produce a
significant valley shift ($+11.6\%$), while an explicit rescheduling authorization adds a further
significant peak reduction ($-8.1\%$). Finally, we show that peak shaving obeys a mechanism bounded by
appliance \emph{substitutability}: a device-targeted intervention shaves the total peak only for
non-substitutable devices (air-conditioning) and not for substitutable ones (induction cooking, which the
agent relocates to the oven and microwave). An adversarial self-audit withdraws one previously claimed
peak-shaving result that does not survive a shared-timeline replication, shows that all headline
$p$-values survive Benjamini--Hochberg correction, and classifies each finding as robust, prompt-sensitive
or withdrawn. We conclude by arguing that the calibrated \emph{capability boundary} of LLM-driven energy
simulation --- not any single effect size --- is the contribution that makes it trustworthy as a
policy-experiment instrument.
\end{abstract}

\noindent\textbf{Keywords:} LLM agents; residential load simulation; demand response; natural-language
intervention; emergent behaviour; peak shaving; validity.

%======================================================================
\section{Introduction}

\subsection{Motivation}
Residential electricity consumption is becoming both more variable and more consequential. Rooftop
photovoltaics, batteries, electric vehicles and heat pumps are turning passive customers into active
prosumers, while demand-side flexibility is increasingly relied upon to absorb the variability of
renewable generation. Traditional load forecasting is ill-equipped for this world: top-down statistical
models treat households as a single aggregate and cannot represent the behavioural mechanisms through
which prices, incentives or social signals actually change consumption. Bottom-up models
\cite{capasso1994,widen2010,richardson2010} reconstruct load from household composition and activity
patterns, but typically reduce behaviour to finite states or Markov chains with fixed transition
probabilities. Such representations struggle with exactly the signals that modern demand-response
programmes use: a dynamic tariff that must be understood and weighed, a social-comparison message, or the
news of an impending heatwave.

Large language model (LLM) generative agents offer a fundamentally more expressive behavioural engine.
The ``generative agents'' architecture of Park et al.\ \cite{park2023} demonstrated that LLM agents
equipped with memory, reflection and planning can produce plausible, long-horizon behaviour, and
subsequent work has scaled such agents to large social simulations \cite{piao2025} and reduced their cost
\cite{koley2025}. In the energy domain, LLM agents have been used to infer single-household consumption
\cite{chetty2025}, to synthesise private household data \cite{almashor2024}, and to parse natural-language
instructions for home energy management \cite{michelon2025}. What has \emph{not} been established is
whether a population of heterogeneous LLM households can be used as a \emph{policy experiment}: whether,
when a tariff, a norm or a disruption is injected purely as text, the population produces aggregate
behaviour that is consistent with the empirical demand-response literature --- and, equally important,
where it is \emph{not}.

\subsection{Research questions}
This paper is organised around four questions. \textbf{RQ1} asks how a large population of heterogeneous
household agents can be generated and mapped to realistic, diverse, minute-resolution load; in this work
we exercise the population machinery at a deliberately modest scale and leave distributional
population-scale validation to future work. \textbf{RQ2} asks how agents respond, at the
appliance-decision level, when prices, taxes, subsidies, norms and news events are injected as natural
language. \textbf{RQ3} asks whether the aggregate responses align with empirical benchmarks and whether
phenomena beyond individual behaviour emerge. \textbf{RQ4} concerns scalability and practicality; it is
explicitly out of scope here, and we return to it only in the discussion as a direction for future work.

\subsection{Contributions}
Our contributions are fourfold. First, we build an end-to-end pipeline that spans
\emph{behaviour planning $\rightarrow$ appliance-level decision-making $\rightarrow$ minute-resolution
load} on top of LLM agents, with all interventions delivered as natural language. Second, we design and
execute a large, controlled experimental programme across three synthetic communities, obeying four
validity rules that we derived from concrete failures, and we report \textbf{seven statistically
significant cross-world event results} together with a validated \textbf{price-response capability} and a
set of well-characterised null results. Third, we provide a mechanistic account of peak shaving: a
device-targeted intervention shaves the total peak only when (i) it is anchored jointly to a
\emph{device} and a \emph{window}, and (ii) the device is not substitutable; this yields the organising
observation that energy saving, peak shaving and load shifting are three distinct behaviours that a single
intervention may or may not produce. Fourth, and perhaps most importantly for the credibility of
LLM-agent energy simulation, we perform an adversarial self-audit, withdraw a previously claimed result
that does not survive a shared-timeline replication, and articulate a calibrated capability boundary:
what such a platform can, cannot, and should not claim.

%======================================================================
\section{Related work}

\subsection{From top-down forecasting to bottom-up behavioural models}
Bottom-up residential load models reconstruct demand from household and appliance inventories and from
occupant activity patterns. Capasso et al.\ \cite{capasso1994} combined socio-economic characteristics,
Monte-Carlo sampling of appliance ownership and probabilistic usage functions to reproduce distribution
level demand without historical aggregate data. Wid\'en and W\"ackelg\r{a}rd \cite{widen2010} derived
activity sequences from time-use surveys via non-homogeneous Markov chains, and Richardson et al.\
\cite{richardson2010} jointly modelled occupancy and appliance usage to produce one-minute-resolution
profiles that matched both distributions and temporal structure of measured data. These models remain the
reference for synthetic load generation, but their behavioural representations are rigid: behaviour is
reduced to a small state space and fixed transition probabilities, and households are largely independent,
so social contagion and coordinated responses to textual or complex signals are out of reach. Agent-based
hybrids such as Xu et al.\ \cite{xu2026} add signal--demand feedback, but their agents remain rule-based.

\subsection{LLM-based generative agents and reasoning}
The generative-agent architecture of Park et al.\ \cite{park2023} --- a memory stream, retrieval,
reflection and planning --- established a general template for behaviour simulation. Scaling it raises
engineering questions of cost and coherence: SALM \cite{koley2025} reduced token consumption through
hierarchical prompting and sustained thousands of steps with an attention-based memory, while
AgentSociety \cite{piao2025} scaled to more than ten thousand urban agents and demonstrated system-level
emergent patterns. At the population level, Chen et al.\ \cite{chen2025} studied consensus-seeking among
LLM agents and showed that agent number, personality and network topology govern convergence. The
behavioural coherence of such agents ultimately rests on reasoning techniques: chain-of-thought prompting
\cite{wei2022} decomposes goals into ordered steps; ReAct \cite{yao2023a} interleaves reasoning with
tool/environment actions; Tree-of-Thoughts \cite{yao2023b} searches over alternative reasoning paths; and
self-consistency \cite{wang2023} samples multiple paths and majority-votes. These techniques are the
methodological substrate that lets an agent plan a day and then decide, segment by segment, which
appliances to run.

In the energy domain, Chetty et al.\ \cite{chetty2025} integrated chain-of-thought reasoning with a
reflective memory module to infer single-household consumption from demographics, reporting a mean
absolute error below 10\% on the Pecan Street dataset; Almashor et al.\ \cite{almashor2024} used private
LLMs to synthesise household energy datasets without sensitive data; and Michelon et al.\
\cite{michelon2025} translated non-standardised natural-language HEMS requests into structured parameters
with high accuracy. These works establish feasibility but remain single-household or small-scale, and none
provides a population-level intervention testbed with empirical alignment. That is the gap this paper
addresses.

\subsection{Empirical benchmarks for alignment}
The empirical literature provides quantitative targets against which a simulated population can be
judged. \emph{Pricing and demand response:} Faruqui and Sergici's survey of fifteen dynamic-pricing
experiments \cite{faruqui2010} reports that plain time-of-use pricing cuts peak demand by roughly
3--6\%, while critical-peak pricing reaches 13--20\%; demand-charge management prototypes cut peak demand
by 10--20\% \cite{escarrega2025}; and Wang et al.\ \cite{wang2021}, using provincial panel data and
questionnaire responses, argue that habit rather than price is the dominant driver of household
consumption. \emph{Non-price behavioural interventions:} social-comparison feedback saves on the order of
1--3\%, with the strongest response among high-usage households \cite{allcott2011,ayres2013}; loss framing
saves about five percentage points more than gain framing \cite{ghesla2019}; monetary incentives can crowd
out intrinsic, social motives \cite{pellerano2017}; the effects of nudges and rebates are strongly
heterogeneous across households \cite{murakami2022}; and the same intervention can work in opposite
directions across ideological groups \cite{costa2010}. \emph{Information:} real-time in-home feedback can
triple the price elasticity of demand \cite{jessoe2012}, and personalised feedback outperforms aggregate
billing \cite{monacchi2015}. \emph{Structural factors:} household-size economies imply that per-capita
demand rises as household size falls \cite{schroder2013}; active and passive peer effects shape the
spatial diffusion of solar panels \cite{barnes2022}; and comfort perception is the main obstacle to
heating sobriety \cite{cabezas2025}. \emph{Disruptions:} empirically grounding LLM agents improves the
realism of behaviour during disruptions \cite{xia2026}, and counterfactual simulation provides a control
paradigm for such events \cite{fidone2025}. \emph{Platforms and methods:} CoRenew \cite{zhang2026} and
Eco3S \cite{wei2026} are agent-based policy-simulation platforms; Gunkel et al.\ \cite{gunkel2023}
analyse the redistributional effects of a uniform electricity tax; and Zhou et al.\ \cite{zhou2016} show
that users whose consumption is more variable are more likely to reduce demand during demand-response
events.

\subsection{Positioning}
Our framework sits at the intersection of these strands. Like bottom-up models, it produces
minute-resolution appliance-level load; like LLM generative agents, it plans behaviour in natural
language; and unlike prior LLM energy work, it is explicitly constructed and stress-tested as a
population-level \emph{policy-intervention instrument} whose outputs are compared, qualitatively, against
the empirical benchmarks above --- with an emphasis on honesty about where the comparison fails.

%======================================================================
\section{The \textsc{LLMWorld} framework}

\subsection{Architecture}
A single entry point, \texttt{run.py}, exposes two modes. The \emph{world} mode synthesises a community
in four stages: household types ($s_1$) $\rightarrow$ persona alignment ($s_2$) $\rightarrow$ household
construction ($s_3$) $\rightarrow$ world assembly ($s_4$). The \emph{simulate} mode generates behaviour,
also in four stages: macro planning ($s_1$) $\rightarrow$ inter-member coordination ($s_2$) $\rightarrow$
contextual enrichment ($s_3$) $\rightarrow$ appliance decision-making ($s_4$). Every intermediate artefact
is persisted as structured JSON, so that any experiment can be reconstructed and re-analysed offline
without re-querying the model. This design choice --- storing the full behavioural trace --- underpins the
audit in Section~\ref{sec:audit}.

\subsection{Domain model}
A household is modelled as a composite of \emph{members}, \emph{rooms} and \emph{appliances}. Appliances
are abstracted into four behavioural classes: \emph{on-demand} (power is drawn only while used, e.g.\ a
lamp or an air-conditioner), \emph{charging} (storage--discharge, e.g.\ a phone or an electric vehicle),
\emph{cycle} (a fixed program with fixed energy per run, e.g.\ a washing machine, dishwasher or oven), and
\emph{always-on} (continuous baseload, e.g.\ a refrigerator). A shared catalog supplies, for each of the
supported appliance types, a rated power and plausible bounds, always-on daily energy, and experiment
metadata --- crucially a \texttt{flexible} flag, a duty cycle, and a season tag --- as well as battery
capacity and state of charge for charging appliances. Charging appliances are additionally bounded by a
battery deficit, so that a device cannot absorb more than one full battery per day even if the agent
requests it.

\subsection{The behaviour-to-load pipeline}
\textbf{Macro planning ($s_1$).} Each member receives a structured profile (age, occupation, habits,
personality) and the day's context (date, weather, active news) and produces a complete 00:00--24:00
timeline of activity segments. The output is constrained to real rooms, must cover the whole day, and
segments must be adjacent; a cross-day carry-over mechanism injects the previous day's end state so that
an agent does not ``teleport home'' at midnight.

\textbf{Coordination ($s_2$).} Where a household has several members, their timelines are reconciled:
shared rooms and resources are sequenced so that, for example, two members do not occupy the kitchen
simultaneously.

\textbf{Enrichment ($s_3$).} Each activity is enriched with the detail required to decide appliance usage
--- for instance, ``cook dinner'' becomes a specific set of cooking steps.

\textbf{Appliance decision ($s_4$).} This is the stage at which interventions bite. For every time segment
of the enriched timeline, the agent decides which appliances to operate, choosing among the actions valid
for each appliance class. The output is validated against the household registry: unknown identifiers are
repaired by appliance family, and operations on room appliances are dropped when the member is out of the
house.

\textbf{Energy calculation.} A separate physical model converts the appliance decisions into a
minute-resolution household load profile (1440 values) and a per-appliance energy breakdown. Cycle
appliances consume their fixed cycle energy, on-demand appliances their rated power weighted by duty
cycle, and idle appliances their standby draw.

\subsection{Intervention surface}
All interventions are injected \emph{as natural language} into the prompts; no behavioural rule is
hard-coded. Pricing policies render tariff text into the decision prompt and include time-of-use pricing
(with a soft, facts-only variant), critical-peak pricing, critical-peak rebates, a uniform electricity
tax, off-peak subsidies, a demand charge, and an EV-delay incentive; social interventions include a fixed
social-norm comparison, its soft variant, a loss-framed norm, and a dynamic peer comparison built from the
previous day's community mean; behaviour-guidance interventions include night setback and a real-time
in-home display. Policies can be combined (e.g.\ \texttt{tou,nudge}) and scheduled over a date window, so
that announcement, effectiveness and removal can be studied. News/events comprise thirteen preset
templates (heatwave, cold-snap, storm, price-rise, supply-margin warning, air-conditioner peak tax,
rebate, rolling-blackout warning, solar subsidy, lockdown, public holiday, work-from-home day, and
transport strike) plus arbitrary custom events and community notices. Preset environment events
(heatwave, cold-snap) additionally adjust the structured weather fields so that the injected text and the
agent's environment agree.

\subsection{Concrete cost context}
Individually, none of the pricing policies above has any effect (Section~\ref{sec:price}); the agents
possess no economic model with which to weigh a tariff. We therefore introduce an optional
\emph{cost-context}: given the active tariff and the household's flexible appliances, the decision prompt
receives a concrete money table. The following is an actual injected example:
\begin{verbatim}
Flexible-appliance costs under today's tariff (peak 16:00-21:00 @0.90; off-peak 22:00-07:00 @0.18):
- bedroom_1_airconditioner: 1.80 AUD now (peak) vs 0.36 AUD off-peak -- save 1.44 per hour
- kitchen_dishwasher:       0.99 AUD now (peak) vs 0.20 AUD off-peak -- save 0.79 per cycle
- bathroom_waterheater:     2.70 AUD now (peak) vs 0.54 AUD off-peak -- save 2.16 per hour
\end{verbatim}
Two variants are supported. In the \emph{strong} variant an explicit sentence authorizes the agent to
move a flexible appliance to an off-peak segment; in the \emph{soft} variant the same figures are given
with a neutral instruction, providing a prompt-bias control. A demand-charge variant instead renders each
appliance's kW $\times$ rate contribution to the daily maximum, and a multi-day variant injects
yesterday's peak-window cost as feedback.

\subsection{Analysis tooling}
A dedicated analysis package computes the quantities used throughout the paper: aggregate population
profiles, per-appliance usage, behaviour--load consistency checks, grouped responses (by energy awareness
or conscientiousness), event response, policy trade-offs, cross-world comparison, load-shape clustering,
forecastability, and a peak decomposition that isolates a single appliance family's contribution to the
peak window. All outputs are persisted as JSON/CSV. An offline suite of 298 unit tests (with the model
mocked) protects the pure logic.

%======================================================================
\section{Experimental design and validity controls}
\label{sec:design}

\subsection{Worlds and households}
We use three synthetic communities with deliberately different structures: \textsc{W1}
(\texttt{world\_838587}, two households), \textsc{W2} (\texttt{world\_172148}, three households) and
\textsc{W3} (\texttt{world\_143345}, three households, including an electric vehicle). Household sizes
range from one to six members, and the households differ in their appliance inventories and in where
their evening peak originates. Unless stated otherwise, an experiment targets one household on one
simulated day and is repeated many times.

\subsection{Design and statistics}
Each comparison is a paired design in which the baseline and the treatment are interleaved so that they
share the same session-time context (``same-era''). Where the intervention acts at the appliance-decision
stage --- i.e.\ pricing and cost-context --- we run only the $s_4$ stage against a \emph{shared}
$s_1$--$s_3$ timeline, which removes planning-level noise and isolates the decision-level effect. Where
the intervention changes \emph{what the household does} --- i.e.\ structural events --- we must run the
full pipeline, because the relevant behaviour is decided during planning. Sample sizes are $n=15$ per arm
unless noted; effect sizes are reported as paired differences with paired $t$ tests, and \$p$-values are
two-sided.

\subsection{Four validity rules}
A substantial part of the engineering effort went into discovering --- and then enforcing --- four rules
without which the results are not interpretable.
\begin{enumerate}
\item \textbf{Drift.} The continuous behaviour of an LLM agent drifts systematically across session time;
two independently sampled runs can differ by more than a treatment effect. We therefore always compare
against a \emph{same-era} baseline, and, wherever possible, prefer \emph{binary or structural} signals
(appliance on/off; out-of-home yes/no) which are immune to drift. The discovery of this artefact led us to
withdraw two earlier ``peak-shaving'' results.
\item \textbf{Power.} Continuous magnitudes are unreliable below $n\approx15$: in a time-of-use
experiment, the estimated peak change moved from $-17.5\%$ at $n=3$ to $+13.1\%$ at $n=9$ to $-3.2\%$
(n.s.) at $n=15$. The run-to-run noise floor is a standard deviation of about $1.0$\,kWh
($\mathrm{CV}\approx12\%$), so literature-scale effects of 3--6\% lie below the resolution of the
instrument.
\item \textbf{Cross-world.} A mechanistic claim is accepted only if it replicates across worlds; several
initially promising effects turned out to be specific to a single household.
\item \textbf{Content, not channel.} An event delivered through the preset template and the same event
delivered as a custom string --- with byte-identical text --- differ on no metric, so the operative
variable is the \emph{content} of the intervention, not the injection path. This resolved a confound that
had confused two earlier experiments.
\end{enumerate}

%======================================================================
\section{Results}
\label{sec:results}

\subsection{Environmental and disruption events}
Table~\ref{tab:events} summarises the disruption events. A heatwave warning causes the air-conditioner to
be switched on where it was previously off (baseline $0/9$ vs heatwave $11/12$, Fisher two-sided
$p\approx3\times10^{-5}$). A lockdown removes almost all out-of-home activity (out-of-home member-days
$13/13\rightarrow1/13$, $p\approx3\times10^{-6}$) and raises daytime (09:00--17:00) energy by $220.7\%$ in
a same-era design; the magnitude is implausibly large in absolute terms and we therefore claim only its
direction. A cold-snap event switches heating from off to on ($0/10$ vs $6/10$, $p\approx0.011$), with a
weaker activation probability than the heatwave. These results are consistent with the disruption-robustness
argument of Xia et al.\ \cite{xia2026}.

\begin{table}[h]
\centering
\small
\begin{tabular}{lll}
\toprule
\textbf{Event} & \textbf{Evidence} & \textbf{Significance} \\
\midrule
Heatwave $\rightarrow$ cooling & baseline AC $0/9$ vs $11/12$ & $p\approx3\times10^{-5}$ \\
Lockdown $\rightarrow$ stay-home & out-of-home $13/13\rightarrow1/13$; daytime $+220.7\%$ & $p\approx3\times10^{-6}$ \\
Cold-snap $\rightarrow$ heating & $0/10$ vs $6/10$ & $p\approx0.011$ \\
\bottomrule
\end{tabular}
\caption{Environmental / disruption events (cross-world).}
\label{tab:events}
\end{table}

\subsection{The stay-home family}
A second, previously unreported group of results concerns \emph{structural stay-home} shocks. A public
holiday raises total daily energy in all seven households tested, by between $19\%$ and $65\%$
($p\le0.019$); a work-from-home day raises it by $30.6\%$, $28.7\%$ and $30.6\%$ in the three worlds
($p\le0.014$, the most consistent effect we observe); and a transport strike raises it by $33.9\%$,
$22.2\%$ and $29.3\%$ ($p\le0.039$). Table~\ref{tab:stayhome} collects these. The common mechanism is
\emph{time at home}: the magnitude of the effect scales with how much the household is normally away
during the day (holiday effect vs baseline daytime share, $r=-0.47$; the weekend-versus-weekday effect
correlates with the holiday effect across households, $r=0.82$). The evening-peak direction, by contrast,
is household-dependent and is not asserted.

\begin{table}[h]
\centering
\small
\begin{tabular}{lll}
\toprule
\textbf{Event} & \textbf{Total-energy change (W1 / W2 / W3)} & \textbf{Significance} \\
\midrule
Public holiday & $+19\%$\ldots$+65\%$ (7 households) & $p\le0.019$ (7/7) \\
Work-from-home & $+30.6\%$ / $+28.7\%$ / $+30.6\%$ & $p\le0.014$ \\
Transport strike & $+33.9\%$ / $+22.2\%$ / $+29.3\%$ & $p\le0.039$ \\
\bottomrule
\end{tabular}
\caption{The stay-home family: structural shocks that raise total energy by increasing time at home.}
\label{tab:stayhome}
\end{table}

\subsection{Price response}
\label{sec:price}
As plain text, none of the pricing policies has a resolvable effect: time-of-use pricing, for example,
changes the evening peak by $-3.2\%$ ($n=15$, n.s.), and a critical-peak price does not shave the peak.
The picture changes decisively once the cost context is injected. Relative to a no-policy baseline, on a
shared timeline, time-of-use pricing with the cost context shifts the peak by $-13.2\%$, $-5.0\%$ and
$-3.5\%$ across the three worlds and raises the valley by $+25.6\%$, $+8.2\%$ and $+8.8\%$
($p\le0.035$); critical-peak pricing with the cost context gives a peak change of $-12.2\%$ and a valley
change of $+23.2\%$; and the rebate variant (critical-peak rebate) gives $-13.7\%$ and $+27.9\%$. Total
energy is essentially unchanged in all cases: this is load shifting, not conservation.

Crucially, we decompose the effect. When the money table is delivered \emph{without} the authorizing
sentence (the soft variant), the valley still rises significantly by $+11.6\%$ ($p=0.019$) while the peak
falls only non-significantly by $-5.5\%$; adding the explicit authorization then produces a further
significant peak reduction of $-8.1\%$ ($p=0.031$) and a valley rise of $+12.5\%$ ($p=0.028$). The
interpretation is therefore two-part: \emph{concrete cost information is itself effective}, and an
actionable permission strengthens the peak effect --- but neither is a clean economic elasticity. The
effect is also dose-dependent: a larger peak/valley gap increases shifting (peak $-7.3\%$, valley
$+14.8\%$, $p\le0.014$), and households whose personas are more price-sensitive shift significantly more
(peak $-7.0\%$, valley $+14.6\%$, $p=0.014$). Reward and penalty framing are statistically
indistinguishable (critical-peak rebate $\approx$ critical-peak tax). On the negative side, a demand charge
remains null even with concrete figures --- it requires a \emph{global} ``lower today's maximum'' action
that the agents do not perform --- a ``free electricity'' window is nominal but non-significant, and
multi-day bill feedback is a clean null (there is no cross-day memory).

\subsection{The peak-shaving mechanism}
The evening peak is dominated by cooking in most households (the induction cooker accounts for
27--54\% of the peak-window energy), and this determines what can be shaved. Targeting a device shaves the
\emph{total} peak only if the device is not substitutable. Targeting the air-conditioner --- which has no
substitute --- reduces air-conditioner energy by $58\%$; targeting the induction cooker reduces its own
load by $99.7\%$, but the agent simply relocates the activity to the oven and microwave, leaving the total
peak unchanged ($-0.2\%$). A device-targeted intervention therefore requires a joint \emph{device}
$\times$ \emph{window} anchor, and generic window-only signals do not shave. This yields the organising
observation that \textbf{energy saving, peak shaving and load shifting are three distinct behaviours}: a
request can cut air-conditioner use while leaving the peak unmoved (saving without shaving); a
critical-peak price can shift load without saving (shifting without saving); and only a device-and-window
intervention on a non-substitutable device produces true peak shaving.

\subsection{Behavioural and social interventions}
Consistent with the pricing nulls, the non-price behavioural interventions are weak or unstable. Social
comparison is directionally inconsistent (it sometimes increases consumption), and a loss-framed norm ---
which is what the literature would predict to be stronger --- produces a counterexample. In-home display
and night setback are null. When framing is layered on top of a concrete price, it ceases to matter:
adding a social-norm message or a loss-framed message to the time-of-use-plus-cost-context intervention
changes the peak by only $+1.4\%$ and $+4.7\%$ respectively, both non-significant. Two combinations are
informative: combining time-of-use pricing with a social norm is non-additive (measured $-13.9\%$ versus
an additive prediction of $-2.3\%$), and combining the air-conditioner peak tax with a generic rebate
\emph{reverses} the tax's effect (the tax alone reduces air-conditioner energy by $58\%$, but adding a
rebate brings it back up by $80\%$), consistent with the crowding-out documented by Pellerano et al.\
\cite{pellerano2017}.

\subsection{Descriptive analyses}
Several analyses characterise the synthetic population. Households differ widely in regularity and split
into recognisable load-shape archetypes (evening-peak, morning-peak and late-night types). Per-capita
daily energy falls as household size rises, qualitatively consistent with the household-size economies of
Schr\"oder et al.\ \cite{schroder2013}. Weekend energy exceeds weekday energy by $3\%$ to $45\%$
depending on the household --- the periodic version of the stay-home effect. A random-forest one-hour-ahead
forecaster improves on a persistence baseline by $31.5\%$ to $57.1\%$, indicating that the generated load
is structurally regular rather than noise. Finally, households whose baseline consumption is more variable
tend to respond more strongly to pricing; the correlation is in the direction reported by Zhou et al.\
\cite{zhou2016} but is not statistically significant at the five households available ($r=0.51$).

%======================================================================
\section{Discussion}
\label{sec:discussion}

\subsection{Two classes of behaviour}
The clearest reading of our results is that the platform contains two qualitatively different kinds of
behaviour. Structural, contextual shocks --- weather and stay-at-home events --- change the \emph{situation}
the agent is reasoning about, and the agent responds through ordinary commonsense reasoning; these effects
are large, robust and cross-world. Generic prices, by contrast, demand an economic trade-off that the
agents do not perform: a tariff is, to them, a sentence, not an incentive. They respond only when the
trade-off is made \emph{concrete} (a per-appliance money figure) and, for peak shaving, \emph{actionable}
(a device and a window). It is tempting to describe this as ``the agents understand prices''; the more
accurate description is that they act on concrete, actionable instructions, and that supplying the money
figures is what makes the instruction concrete.

\subsection{Alignment with the empirical literature}
The directions of our results align with the empirical record. Peak shifting under concrete pricing is
consistent with Faruqui and Sergici \cite{faruqui2010}; the disruption results are consistent with Xia et
al.\ \cite{xia2026}; the non-additivity of combinations and the crowding-out of a device tax by a rebate
mirror Pellerano et al.\ \cite{pellerano2017}; the fragility of framing effects without concrete anchors
echoes the habit-dominance argument of Wang et al.\ \cite{wang2021}. The \emph{magnitudes}, however,
systematically exceed the empirical values, which we attribute to LLM over-compliance with explicit
natural-language directives: an instruction-following agent will overshoot relative to a habit-bound human
household. For this reason we treat alignment as qualitative and decline to claim absolute magnitudes.

\subsection{A calibrated capability boundary}
We therefore state the boundary explicitly. The platform \emph{can} reliably give (i) the direction of
structural weather and stay-home events, (ii) peak shifting and valley filling under concrete pricing, and
(iii) the device-level response of non-substitutable appliances. It \emph{cannot} give literature-scale
(3--6\%) text-price elasticity, global (demand-charge) shaving, cross-day learning or habit, and it cannot
represent photovoltaic-adoption or electric-vehicle valley phenomena because those mechanisms are absent.
Naming this boundary is not a concession; it is what turns an impressive demo into a usable instrument,
because it tells a policy analyst in advance which questions the tool can answer.

\subsection{Methodological contributions}
Three methodological findings may outlast the platform itself. First, continuous LLM-agent behaviour
drifts across session time, so same-era baselines or binary signals are mandatory --- a result that
invalidated two of our own earlier conclusions. Second, magnitudes below $n\approx15$ are unreliable and
should not be reported. Third, the effect of an intervention depends on its \emph{content}, not on the
channel through which it is delivered. Together with the four rules of Section~\ref{sec:design}, these
findings form a validity protocol that we recommend for LLM-agent simulation studies generally.

%======================================================================
\section{Adversarial validity audit}
\label{sec:audit}
Because the experimental programme is large, we subjected every conclusion to a deliberate attempt at
disconfirmation, and we report the outcome even where it is unflattering.

\noindent\textbf{Robust.} The seven event results and the null set are robust in direction. To guard
against multiple-comparison artefacts, we applied the Benjamini--Hochberg procedure across the eight
primary tests; all eight survive (the largest adjusted $p$-value is $0.039$).

\noindent\textbf{Prompt-sensitive (preliminary).} The price-shaving result is partly instruction-driven
and should not be reported as a pure elasticity (Section~\ref{sec:price}). The heatwave result is not
purely language-only, because the preset event also adjusts a structured weather field. The season result
largely follows the prompt's own season rule and is therefore weak evidence of emergent seasonal
behaviour.

\noindent\textbf{Withdrawn or demoted.} One previously claimed result --- that an air-conditioner peak tax
reduces the total evening peak by $22.5\%$ --- is \emph{withdrawn}: on a shared timeline the corresponding
comparison gives $+6.1\%$ (n.s.), and only the device-level air-conditioner reduction of $58\%$ survives.
Two further results are demoted to weak evidence: the non-additivity of the air-conditioner tax and rebate
is based on $n=9$, and the substitution-based explanation of the cooking peak did not replicate on a
full five-member household.

\noindent\textbf{Cross-cutting risks.} We also record risks that we could only partially mitigate:
over-compliance inflating magnitudes; interventions that touch structured fields rather than language
alone; a household-specific ``generic-event reactivity'' that inflates load on one community and confounds
several comparisons there; the post-hoc tuning of the authorization wording in the cost context; the
single-member scope of most experiments; and platform-version drift during the study, as analysis and
weather modules were extended mid-programme.

%======================================================================
\section{Limitations}
The weather model is a stub beyond the date-derived season and the event overrides, and the seasonal
temperatures are nominal. We did not align against public occupant-driven stochastic load data, so the
alignment with the empirical literature is qualitative rather than distributional. The population is
exercised at no more than five households per world, and the two research questions that concern scale and
cost (RQ1-at-scale and RQ4) are out of scope. Most experiments simulate a single household member, and
multi-member validation was performed only for a subset of effects. Absolute magnitudes are inflated by
over-compliance and are not claimed. Finally, several phenomena of practical interest --- photovoltaic
adoption, electric-vehicle valley filling, and any form of cross-day learning --- cannot currently be
studied on this platform.

%======================================================================
\section{Conclusion}
This paper set out to ask whether a population of heterogeneous LLM household agents can be used as an
instrument for natural-language energy-policy experiments, and to answer that question honestly. The
instrument we built, \textsc{LLMWorld}, produces emergent, minute-resolution behaviour across an
end-to-end behaviour-to-load pipeline, and it yields seven statistically significant, cross-world event
results, a validated price-response capability, and a mechanistic account of peak shaving governed by a
dose $\times$ substitutability $\times$ device/window rule. Equally important, it yields a precise
negative space: text-only prices, global demand charges, cross-day learning, and photovoltaic and electric
vehicle effects are not attainable, and one previously claimed result is withdrawn. We argue that this
calibrated boundary --- the explicit statement of what an LLM-agent energy simulator can and cannot
establish --- is the contribution that ultimately makes such platforms credible, and we release the
pipeline, the intervention surface and the full experimental record to support replication.

%======================================================================
\begin{thebibliography}{99}
\small
\bibitem{albadi2008} Albadi, M. H., \& El-Saadany, E. F. (2008). A summary of demand response in electricity markets. \emph{Electric Power Systems Research}, 78(11), 1989--1996.
\bibitem{alexeenko2023} Alexeenko, P., \& Bitar, E. (2023). Achieving reliable coordination of residential plug-in electric vehicle charging: A pilot study. arXiv:2112.04559.
\bibitem{allcott2011} Allcott, H. (2011). Social norms and energy conservation. \emph{Journal of Public Economics}, 95(9--10), 1082--1095.
\bibitem{almashor2024} Almashor, M., et al. (2024). Can private LLM agents synthesize household energy consumption data? \emph{Energy and AI}, 16, 100360.
\bibitem{ayres2013} Ayres, I., Raseman, S., \& Shih, A. (2013). Evidence from two large field experiments that peer comparison feedback can reduce residential energy usage. \emph{Journal of Law, Economics, and Organization}, 29(5), 992--1022.
\bibitem{barnes2022} Barnes, J., Islam, M. R., \& Morshed, S. (2022). Passive and active peer effects in the spatial diffusion of residential solar panels. \emph{Energy Research \& Social Science}, 86, 102458.
\bibitem{cabezas2025} Cabezas-Rivi\`ere, N., et al. (2025). Identifying and understanding obstacles to heating sobriety and thermal comfort in collective housing. arXiv:2512.16949.
\bibitem{capasso1994} Capasso, A., Grattieri, W., Lamedica, R., \& Prudenzi, A. (1994). A bottom-up approach to residential load modeling. \emph{IEEE Transactions on Power Systems}, 9(2), 957--964.
\bibitem{chen2025} Chen, H., et al. (2025). Multi-agent consensus seeking via large language models. arXiv:2310.20151.
\bibitem{chetty2025} Chetty, N., et al. (2025). An LLM framework for inferring household energy consumption through behaviour simulation. \emph{Applied Energy}, 355, 122335.
\bibitem{costa2010} Costa, D. L., \& Kahn, M. E. (2010). Energy conservation ``nudges'' and environmentalist ideology. \emph{Journal of the European Economic Association}, 11(3), 680--702.
\bibitem{escarrega2025} Escarrega, C., et al. (2025). Demand charge management: Prototype design and testing. arXiv:2509.10713.
\bibitem{faruqui2010} Faruqui, A., \& Sergici, S. (2010). Household response to dynamic pricing of electricity: A survey of 15 experiments. \emph{Journal of Regulatory Economics}, 38(2), 193--225.
\bibitem{fidone2025} Fidone, G., et al. (2025). Evaluating online moderation via LLM-powered counterfactual simulations. arXiv:2511.07204.
\bibitem{ghesla2019} Ghesla, C., Grieder, M., \& Schmitz, J. (2019). Pro-environmental incentives and loss aversion: A field experiment on electricity saving behavior. \emph{Energy Policy}, 128, 102--112.
\bibitem{gunkel2023} Gunkel, P. A., et al. (2023). Uniform taxation of electricity: Incentives for flexibility and cost redistribution among household categories. arXiv:2306.11566.
\bibitem{jessoe2012} Jessoe, K., \& Rapson, D. (2012). Knowledge is (less) power: Experimental evidence from residential energy use. \emph{American Economic Review}, 104(4), 1417--1438.
\bibitem{jin2021} Jin, L., et al. (2021). Investigating underlying drivers of variability in residential energy usage patterns with daily load shape clustering of smart meter data. arXiv:2102.11027.
\bibitem{koley2025} Koley, S. (2025). SALM: A multi-agent framework for language model-driven social network simulation. arXiv:2505.09081.
\bibitem{michelon2025} Michelon, F., Zhou, Y., \& Morstyn, T. (2025). Large language model interface for home energy management systems. \emph{Applied Energy}, 358, 122590.
\bibitem{monacchi2015} Monacchi, A., et al. (2015). An open solution to provide personalized feedback for building energy management. arXiv:1505.01311.
\bibitem{murakami2022} Murakami, K., Shimada, K., \& Tanaka, M. (2022). Heterogeneous treatment effects of nudge and rebate: Causal machine learning in a field experiment on electricity conservation. \emph{Energy Economics}, 105, 105737.
\bibitem{park2023} Park, J. S., et al. (2023). Generative agents: Interactive simulacra of human behavior. In \emph{Proc.\ UIST '23}.
\bibitem{pellerano2017} Pellerano, J. A., et al. (2017). Do extrinsic incentives undermine social norms? Evidence from a field experiment in energy conservation. \emph{Environmental and Resource Economics}, 67(3), 413--431.
\bibitem{piao2025} Piao, J., et al. (2025). AgentSociety: Large-scale LLM-driven generative agents for social simulation. arXiv.
\bibitem{richardson2010} Richardson, I., Thomson, M., Infield, D., \& Clifford, C. (2010). Domestic electricity use: A high-resolution energy demand model. \emph{Energy and Buildings}, 42(10), 1878--1887.
\bibitem{schroder2013} Schr\"oder, C., et al. (2013). Household formation and residential energy demand: Evidence from Japan. \emph{Energy Economics}.
\bibitem{thorvaldsen2021} Thorvaldsen, K., et al. (2021). Long-term value of flexibility from flexible assets in building operation. arXiv:2105.11952.
\bibitem{wang2021} Wang, Y., et al. (2021). Electricity price and habits: Which would affect household electricity consumption? \emph{Energy Policy}.
\bibitem{wang2023} Wang, X., et al. (2023). Self-consistency improves chain of thought reasoning in language models. In \emph{Proc.\ ICLR}.
\bibitem{wei2022} Wei, J., et al. (2022). Chain-of-thought prompting elicits reasoning in large language models. In \emph{Proc.\ NeurIPS}.
\bibitem{wei2026} Wei, Y., et al. (2026). Eco3S: Complex socio-economic system simulation via agent-based models. arXiv:2607.26588.
\bibitem{widen2010} Wid\'en, J., \& W\"ackelg\r{a}rd, E. (2010). A high-resolution stochastic model of domestic activity patterns and electricity demand. \emph{Applied Energy}, 87(6), 1880--1892.
\bibitem{xia2026} Xia, H., et al. (2026). Empirical grounding improves the realism of LLM agents simulating human behavior during disruptions. arXiv:2607.17437.
\bibitem{xu2026} Xu, et al. (2026). Agent-based modeling and neural network for residential customer demand response.
\bibitem{yao2023a} Yao, S., et al. (2023a). ReAct: Synergizing reasoning and acting in language models. In \emph{Proc.\ ICLR}.
\bibitem{yao2023b} Yao, S., et al. (2023b). Tree of thoughts: Deliberate problem solving with large language models. In \emph{Proc.\ NeurIPS}.
\bibitem{zhang2026} Zhang, et al. (2026). CoRenew: A large language model agent-based policy simulation platform for multifamily residential redevelopment. arXiv:2607.25447.
\bibitem{zhou2016} Zhou, D., Balandat, M., \& Tomlin, C. (2016). Residential demand response targeting using machine learning with observational data. arXiv:1607.00595.
\end{thebibliography}

\appendix
\section{Reproducibility}
The repository contains the simulation entry point, the engine modules for pricing, tariff/cost-context,
news/events and weather, the world- and simulate-stage steps, the analysis package, 298 offline unit tests
with the language model mocked, and the complete experimental record of 222 rounds. Every intervention is
delivered in natural language; no behavioural rule is hard-coded.

\end{document}
